\documentclass[10pt,a4paper]{amsart}
\usepackage{amsmath,amssymb,amsthm,graphicx}
% Retained minimal preamble for the Stacks Project chapter more-algebra.tex (More on Algebra).
% The chapter file pinned-original-source.tex carries only \input{preamble} on its first line, so no
% preamble text lives inside the chapter itself.  The Stacks Project uses one single shared preamble for
% every chapter; the snapshot stages it identically as sources/stacks-chow/preamble.tex and
% sources/stacks-etale-cohomology/preamble.tex.  The block below is copied mechanically, line for line,
% from sources/stacks-chow/preamble.tex lines 14-18 (the Xy-pic setup) and 154-217 (the native theorem
% declarations together with the mathematical macros).  Nothing was added, removed or redefined; the
% wrapper class and the amsmath/amssymb/amsthm/graphicx base packages are supplied by the assembler.
% Chapter: more-algebra.tex -- More on Algebra.
% For commutative diagrams we use Xy-pic
\usepackage[all]{xy}

% We use 2cell for 2-commutative diagrams.
\xyoption{2cell}

% Theorem environments.
%
\theoremstyle{plain}
\newtheorem{theorem}[subsection]{Theorem}
\newtheorem{proposition}[subsection]{Proposition}
\newtheorem{lemma}[subsection]{Lemma}

\theoremstyle{definition}
\newtheorem{definition}[subsection]{Definition}
\newtheorem{example}[subsection]{Example}
\newtheorem{exercise}[subsection]{Exercise}
\newtheorem{situation}[subsection]{Situation}

\theoremstyle{remark}
\newtheorem{remark}[subsection]{Remark}
\newtheorem{remarks}[subsection]{Remarks}

\numberwithin{equation}{subsection}

% Macros
%
\def\lim{\mathop{\mathrm{lim}}\nolimits}
\def\colim{\mathop{\mathrm{colim}}\nolimits}
\def\Spec{\mathop{\mathrm{Spec}}}
\def\Hom{\mathop{\mathrm{Hom}}\nolimits}
\def\Ext{\mathop{\mathrm{Ext}}\nolimits}
\def\SheafHom{\mathop{\mathcal{H}\!\mathit{om}}\nolimits}
\def\SheafExt{\mathop{\mathcal{E}\!\mathit{xt}}\nolimits}
\def\Sch{\mathit{Sch}}
\def\Mor{\mathop{\mathrm{Mor}}\nolimits}
\def\Ob{\mathop{\mathrm{Ob}}\nolimits}
\def\Sh{\mathop{\mathit{Sh}}\nolimits}
\def\NL{\mathop{N\!L}\nolimits}
\def\CH{\mathop{\mathrm{CH}}\nolimits}
\def\proetale{{pro\text{-}\acute{e}tale}}
\def\etale{{\acute{e}tale}}
\def\QCoh{\mathit{QCoh}}
\def\Ker{\mathop{\mathrm{Ker}}}
\def\Im{\mathop{\mathrm{Im}}}
\def\Coker{\mathop{\mathrm{Coker}}}
\def\Coim{\mathop{\mathrm{Coim}}}

% Boxtimes
%
\DeclareMathSymbol{\boxtimes}{\mathbin}{AMSa}{"02}

%
% Macros for moduli stacks/spaces
%
\def\QCohstack{\mathcal{QC}\!\mathit{oh}}
\def\Cohstack{\mathcal{C}\!\mathit{oh}}
\def\Spacesstack{\mathcal{S}\!\mathit{paces}}
\def\Quotfunctor{\mathrm{Quot}}
\def\Hilbfunctor{\mathrm{Hilb}}
\def\Curvesstack{\mathcal{C}\!\mathit{urves}}
\def\Polarizedstack{\mathcal{P}\!\mathit{olarized}}
\def\Complexesstack{\mathcal{C}\!\mathit{omplexes}}
% \Pic is the operator that assigns to X its picard group, usage \Pic(X)
% \Picardstack_{X/B} denotes the Picard stack of X over B
% \Picardfunctor_{X/B} denotes the Picard functor of X over B
\def\Pic{\mathop{\mathrm{Pic}}\nolimits}
\def\Picardstack{\mathcal{P}\!\mathit{ic}}
\def\Picardfunctor{\mathrm{Pic}}
\def\Deformationcategory{\mathcal{D}\!\mathit{ef}}

% --- appended 2026-10-12 (Chief): the project own definitions of the reference/slogan/history
% display environments, copied mechanically from the shared Stacks preamble lines 8-12
% (\usepackage{verbatim} plus the three \newenvironment lines, which the chapter preamble keeps
% verbatim from the project preamble). Nothing added, removed or redefined: these environments
% carry no mathematical content and the project itself typesets them as comments. Without these
% three lines a counted statement whose own text uses them fails with Environment slogan undefined.
\usepackage{verbatim}
\newenvironment{reference}{\comment}{\endcomment}
\newenvironment{slogan}{\comment}{\endcomment}
\newenvironment{history}{\comment}{\endcomment}
\usepackage{cleveref}
\crefname{originaldefinition}{original definition}{original definitions}
\Crefname{originaldefinition}{Original definition}{Original definitions}
\makeatletter
\newlabel{definition-K-flat}{{14185}{source}{}{sourcepointer-1}{}}
\newlabel{definition-K-flat@cref}{{[originaldefinition][0][]14185}{[1][0][]source}{}{}{}}
\makeatother
\crefname{originaldefinition}{original definition}{original definitions}
\Crefname{originaldefinition}{Original definition}{Original definitions}
\makeatletter
\newlabel{definition-perfect}{{18871}{source}{}{sourcepointer-2}{}}
\newlabel{definition-perfect@cref}{{[originaldefinition][0][]18871}{[1][0][]source}{}{}{}}
\makeatother
\crefname{originallemma}{original lemma}{original lemmas}
\Crefname{originallemma}{Original lemma}{Original lemmas}
\makeatletter
\newlabel{derived-lemma-morphisms-from-projective-complex}{{064B}{source}{}{sourcepointer-3}{}}
\newlabel{derived-lemma-morphisms-from-projective-complex@cref}{{[originallemma][0][]064B}{[1][0][]source}{}{}{}}
\makeatother
\crefname{originallemma}{original lemma}{original lemmas}
\Crefname{originallemma}{Original lemma}{Original lemmas}
\makeatletter
\newlabel{lemma-RHom-out-of-projective}{{18708}{source}{}{sourcepointer-4}{}}
\newlabel{lemma-RHom-out-of-projective@cref}{{[originallemma][0][]18708}{[1][0][]source}{}{}{}}
\makeatother
\crefname{originallemma}{original lemma}{original lemmas}
\Crefname{originallemma}{Original lemma}{Original lemmas}
\makeatletter
\newlabel{lemma-derived-tor-homotopy}{{14104}{source}{}{sourcepointer-5}{}}
\newlabel{lemma-derived-tor-homotopy@cref}{{[originallemma][0][]14104}{[1][0][]source}{}{}{}}
\makeatother
\crefname{originallemma}{original lemma}{original lemmas}
\Crefname{originallemma}{Original lemma}{Original lemmas}
\makeatletter
\newlabel{lemma-derived-tor-quasi-isomorphism}{{14284}{source}{}{sourcepointer-6}{}}
\newlabel{lemma-derived-tor-quasi-isomorphism@cref}{{[originallemma][0][]14284}{[1][0][]source}{}{}{}}
\makeatother
\crefname{originallemma}{original lemma}{original lemmas}
\Crefname{originallemma}{Original lemma}{Original lemmas}
\makeatletter
\newlabel{lemma-dual-perfect-complex}{{19176}{source}{}{sourcepointer-7}{}}
\newlabel{lemma-dual-perfect-complex@cref}{{[originallemma][0][]19176}{[1][0][]source}{}{}{}}
\makeatother
\crefname{originallemma}{original lemma}{original lemmas}
\Crefname{originallemma}{Original lemma}{Original lemmas}
\makeatletter
\newlabel{lemma-tensor-product-K-flat}{{14224}{source}{}{sourcepointer-8}{}}
\newlabel{lemma-tensor-product-K-flat@cref}{{[originallemma][0][]14224}{[1][0][]source}{}{}{}}
\makeatother
\title{The Hom complex from a bounded above projective complex representing a perfect object to a K-flat complex is again K-flat}
\author{The Stacks Project Authors}
\date{stacks-project:more-algebra.tex:local-snapshot-20261012; GNU FDL 1.2 or later}
\begin{document}
\maketitle
\setcounter{subsection}{1}
\noindent\textit{Editorial note.} All delivered mathematical text is the author's own, copied from the pinned Stacks Project chapter file more-algebra.tex; no statement or proof step was rewritten, reordered or removed.\par\medskip
\noindent\textit{Editorial note.} The chapter file pinned-original-source.tex carries only \textbackslash{}input\{preamble\} on its first line, so no preamble text lives inside the chapter. The retained file sources/stacks-more-algebra/retained-preamble.tex is a mechanical, line-for-line copy of the project's single shared preamble as staged in the snapshot; nothing was added, removed or redefined and no package was added.\par\medskip
\noindent\textit{Editorial note.} The source URL is the source-level Stacks Project URL, as required when more than one card is produced from one source; the per-claim tag is carried by claim\_key.\par\medskip
\noindent\textit{Editorial note.} The author refers to 'Lemma ... and its proof' when identifying the Hom complex with the total tensor complex; that pointer is recorded as a non-load-bearing navigation reference. The eight references inside the counted proof span are recorded as non-load-bearing navigation pointers, mixing cross-chapter tags with same-chapter label lines.\par\medskip
\noindent\textit{Allowed standard prerequisites (editorial).}\begin{itemize}
\item K-flat complexes and their homotopy invariance under the derived tensor product. used at proof lines 28730-28733 and 28743-28744. Is the property being established and the invariance that permits the replacement by F\textasciicircum{}bullet.
\item Representation of a perfect object by a finite complex of finite projective modules, and homotopy equivalence of complexes of projectives that are isomorphic in the derived category. used at proof lines 28721-28727. Supplies the finite complex F\textasciicircum{}bullet and the homotopy equivalence.
\item The dual complex E\textasciicircum{}bullet of a finite complex of finite projectives and the identification of Hom with a total tensor complex. invoked at proof lines 28734-28742. Rewrites the Hom complex in the form to which the tensor-product K-flatness applies.
\end{itemize}\medskip
\ifdefined\hypertarget\hypertarget{sourcepointer-1}{}\fi
\noindent\textit{Source reference.} Original definition 14185: sources/stacks-more-algebra/pinned-original-source.tex, source lines 14185--14185. Same-chapter navigation pointer; the statement it names is standard for the section and the argument that consumes it stays inside the counted proof span.\par\medskip
\ifdefined\hypertarget\hypertarget{sourcepointer-2}{}\fi
\noindent\textit{Source reference.} Original definition 18871: sources/stacks-more-algebra/pinned-original-source.tex, source lines 18871--18871. Same-chapter navigation pointer; the statement it names is standard for the section and the argument that consumes it stays inside the counted proof span.\par\medskip
\ifdefined\hypertarget\hypertarget{sourcepointer-3}{}\fi
\noindent\textit{Source reference.} Original lemma 064B: focused-seeds/AGP/sources/stacks-snapshot-20261011/tags.txt, source lines 7860--7860. Cross-chapter navigation pointer into the Stacks Project; the standard fact it supplies is applied inside the counted proof span and the argument resting on it stays there.\par\medskip
\ifdefined\hypertarget\hypertarget{sourcepointer-4}{}\fi
\noindent\textit{Source reference.} Original lemma 18708: sources/stacks-more-algebra/pinned-original-source.tex, source lines 18708--18708. Same-chapter navigation pointer; the statement it names is standard for the section and the argument that consumes it stays inside the counted proof span.\par\medskip
\ifdefined\hypertarget\hypertarget{sourcepointer-5}{}\fi
\noindent\textit{Source reference.} Original lemma 14104: sources/stacks-more-algebra/pinned-original-source.tex, source lines 14104--14104. Same-chapter navigation pointer; the statement it names is standard for the section and the argument that consumes it stays inside the counted proof span.\par\medskip
\ifdefined\hypertarget\hypertarget{sourcepointer-6}{}\fi
\noindent\textit{Source reference.} Original lemma 14284: sources/stacks-more-algebra/pinned-original-source.tex, source lines 14284--14284. Same-chapter navigation pointer; the statement it names is standard for the section and the argument that consumes it stays inside the counted proof span.\par\medskip
\ifdefined\hypertarget\hypertarget{sourcepointer-7}{}\fi
\noindent\textit{Source reference.} Original lemma 19176: sources/stacks-more-algebra/pinned-original-source.tex, source lines 19176--19176. Same-chapter navigation pointer; the statement it names is standard for the section and the argument that consumes it stays inside the counted proof span.\par\medskip
\ifdefined\hypertarget\hypertarget{sourcepointer-8}{}\fi
\noindent\textit{Source reference.} Original lemma 14224: sources/stacks-more-algebra/pinned-original-source.tex, source lines 14224--14224. Same-chapter navigation pointer; the statement it names is standard for the section and the argument that consumes it stays inside the counted proof span.\par\medskip
\begin{lemma}
\label{lemma-hom-complex-K-flat}
Let $R$ be a ring. Let $P^\bullet$ be a bounded above complex
of projective $R$-modules. Let $K^\bullet$ be a K-flat complex
of $R$-modules. If $P^\bullet$ is a perfect object of $D(R)$,
then $\Hom^\bullet(P^\bullet, K^\bullet)$ is K-flat and
represents $R\Hom_R(P^\bullet, K^\bullet)$.
\end{lemma}

\begin{proof}
The last statement is Lemma \ref{lemma-RHom-out-of-projective}.
Since $P^\bullet$ represents a perfect object, there exists a
finite complex of finite projective $R$-modules $F^\bullet$
such that $P^\bullet$ and $F^\bullet$ are isomorphic in $D(R)$, see
Definition \ref{definition-perfect}.
Then $P^\bullet$ and $F^\bullet$ are homotopy equivalent, see
Derived Categories, Lemma
\ref{derived-lemma-morphisms-from-projective-complex}.
Then $\Hom^\bullet(P^\bullet, K^\bullet)$
and $\Hom^\bullet(F^\bullet, K^\bullet)$
are homotopy equivalent. Hence the first is K-flat if and
only if the second is (follows from
Definition \ref{definition-K-flat} and
Lemma \ref{lemma-derived-tor-homotopy}).
It is clear that
$$
\Hom^\bullet(F^\bullet, K^\bullet) =
\text{Tot}(E^\bullet \otimes_R K^\bullet)
$$
where $E^\bullet$ is the dual complex to $F^\bullet$
with terms $E^n = \Hom_R(F^{-n}, R)$, see
Lemma \ref{lemma-dual-perfect-complex}
and its proof.
Since $E^\bullet$ is a bounded complex of projectives
we find that it is K-flat by Lemma \ref{lemma-derived-tor-quasi-isomorphism}.
Then we conclude by Lemma \ref{lemma-tensor-product-K-flat}.
\end{proof}

\end{document}
